% Job posting: https://ca.linkedin.com/jobs/view/senior-audio-programmer-at-ubisoft-montr%C3%A9al-4438720654
\documentclass[11pt]{article}
\usepackage[left=1in,top=1in,right=1in,bottom=1in]{geometry}
\usepackage[hidelinks]{hyperref}
\usepackage{setspace}
\usepackage[T1]{fontenc}
\usepackage{newtxtext,newtxmath}
\ifdefined\pdfgentounicode
  \input{glyphtounicode}
  \pdfgentounicode=1
\fi
\setstretch{1.1}
\pagenumbering{gobble}
\begin{document}
\pagestyle{empty}

\begin{flushleft}
Nishal Stanislaus Silva, PhD \\
Toronto, ON, Canada \\
nishal.silva@hotmail.com \,|\, +1 613 200 2030 \\
\href{https://nishal.xyz}{nishal.xyz} \,|\, \href{https://www.linkedin.com/in/nishal-silva}{linkedin.com/in/nishal-silva} \,|\, \href{https://github.com/ns2max}{github.com/ns2max}
\end{flushleft}

\vspace{1em}

Dear Hiring Manager,

I build real-time C++ audio systems under hard latency budgets, and I want to bring that work into a game engine. As a postdoctoral researcher at the University of Trento, I owned the full pipeline for an embedded audio pattern-detection system: 14 ms inference on a Raspberry Pi 4, F1 0.76, running at 31.4\% CPU -- 74x faster than the 1038 ms baseline it replaced -- deployed inside Elk Audio OS as a VST plugin with OSC control. I also maintain Nebula, an open-source C++17 audio-feature-extraction library with per-feature latency benchmarks on ARM and Raspberry Pi hardware, which is direct, public evidence of production-grade C++ audio code shipped and profiled on constrained hardware, the same kind of budget a real-time game audio engine has to live inside.

That technical background sits on top of 15+ years as a performing and session guitarist (Trinity College London Rock \& Pop Guitar, Grade 8), which gives me a musician's ear for audio quality and mix feel, not just the signal-processing math behind it. I have also built and shipped interactive real-time audio-visual systems directly: during my postdoc I co-designed and ran two live multisensory concerts where a real-time audio pipeline was synchronized with mixed-reality visuals, stage lighting, and haptics across four performers simultaneously, evaluated in front of a public audience of 20 with 6 performers on stage. I also spent time as a visiting researcher at McGill University's IDMIL/CIRMMT in Montreal, where I built a self-powered smart electric guitar fusing IMU and audio streams in real time -- so I already have working ties to the city this role is based in.

I have not worked in the games industry before, and I want to be direct about that rather than talk around it. What I bring instead is a combination that is uncommon in either direction: research-grade real-time audio DSP engineering, proven on embedded hardware with published, verified numbers, paired with a working musician's sensibility for what audio actually needs to feel like to a listener and a performer. I see this role as a deliberate pivot into games, not a stretch from adjacent industry experience, and I am confident that the depth transfers directly to the constraints of a AAA audio engine.

I am a Canadian permanent resident, eligible to work in Canada without sponsorship, based in Toronto with existing ties to Montreal through my time at McGill, and available to start immediately. I would welcome the opportunity to discuss how my real-time audio engineering background could contribute to Ubisoft Montreal's audio team.

Sincerely, \\
Nishal Stanislaus Silva

\end{document}
